\paragraph{The finite Schwartz function.}
All local functions below are zero off their specified supports.
Set $V=B^0$ and $Q_B(x)=\operatorname{Nrd}(x)=-x^2$.
For $p\nmid2Nm\ell_0$ use $1_{R_p^0}$.
For odd $p^e\Vert N$, use the split integral unit support.  The
cyclic level quotient selects one of the two eigensummands and its
root $r\in\mathbb Z_p^\times$.  The projectors
$(x+r)/(2r)$ and $(r-x)/(2r)$ are integral, so reduction and Hensel
lifting show that the selected root modulo any fixed power of $p$
is locally constant.  Weight it by a unit character $\alpha_p(r)$
of finite $2$-power order with
\[
 \alpha_p(-1)=w_{p^e}\chi_{h_*}(p)^e. \tag{O.3}
\]
Such a character exists: at an odd prime the cyclic $2$-primary
part of $\mathbb F_p^\times$ admits a character taking $-1$ to
$-1$, and the trivial character supplies the other parity.
Conjugation by $R_p^\times$ preserves the cyclic quotient and its
selected summand.  Multiplication of $x$ by a unit $u$ preserves
the summands and replaces $r$ by $ur$, giving multiplier
$\alpha_p(u)$.

At two require $e_x=(1+x)/2\in R_2$ and
$x^2=r^2\in\mathbb Z_2^{\times2}$.  The eigenvalues
$a=(1+r)/2$ and $a'=(1-r)/2$ of $e_x$ differ by the unit $r$.
Its projectors $(e_x-a')/r$ and $(a-e_x)/r$ are integral and remain
distinct modulo two.  If $2^e\Vert N$, the cyclic quotient is a
free rank-one $\mathbb Z/2^e$-module.  Of these complementary
projectors exactly one acts as the identity on it, since the only
idempotents of $\mathbb Z/2^e$ are zero and one.  This defines the
selected root without a choice of its parity.  For every odd $u$,
\[
 e_{ux}=u e_x+(1-u)/2.
\]
This affine transformation preserves the two projectors
individually, so the selected root changes to $ur$, even when
the two eigenvalue parities switch.  The same root character rule
therefore applies at two.  A character with either prescribed
parity exists, using the trivial character or the character modulo
four.  If $2\nmid N$ there is no root character at two.

For $p\mid m$, require $x\in M_2(\mathbb Z_p)^0$ and $v_pQ_B(x)=1$.
The reduction of $x$ is a nonzero nilpotent: zero reduction would
force $p^2\mid Q_B(x)$.  All nonzero nilpotents are conjugate to
$X_0=\left(\begin{smallmatrix}0&1\\0&0\end{smallmatrix}\right)$.
The centralizer of $X_0$ in $\operatorname{GL}_2(\mathbb F_p)$ is
$\{\left(\begin{smallmatrix}a&b\\0&a\end{smallmatrix}\right):a\ne0\}$
and has square determinant.  Thus
\[
 \epsilon_p(gX_0g^{-1})=(\det g/p)
\]
is well-defined.  Conjugation by $u$ multiplies it by $(\det u/p)$,
and scalar multiplication by $u\in\mathbb Z_p^\times$ multiplies
it by $(u/p)$, as follows from conjugation of $X_0$ by
$\operatorname{diag}(u,1)$.

At $\ell_0$ require $x\in R_{\ell_0}^0$ with nonzero reduction
$\bar x=t\beta$ in the trace-zero line of $\mathbb F_{\ell_0^2}$,
and use $(t/\ell_0)$.  Order units act trivially on the residue
field.  A division uniformizer acts by Frobenius, which sends
$\beta$ to $-\beta$, giving multiplier $-1$.
Unit scaling gives $(u/\ell_0)$.  Moreover $\beta^2$ is a nonsquare
in $\mathbb F_{\ell_0}$: otherwise $\beta=\pm a$ for
$a\in\mathbb F_{\ell_0}$, contrary to its nonzero trace-zero condition.
Both $-1$ and $\beta^2$ are nonsquares, so $Q_B(x)=-t^2\beta^2$
is a nonzero square.  Since $(m/\ell_0)=1$, extracted indices have
$(D/\ell_0)=1$.

These supports are compact and open, and all weights are locally
constant.  Additional fixed unit-squareclass conditions multiply
them by locally constant characteristic functions.  They are
preserved by $x\mapsto ux$, which multiplies the norm by $u^2$.
Let $\mathcal O$ be the valuation ring at a place above two of the
field generated by the fixed root-character values at primes
dividing $N$.  All other weights are quadratic, so this field is
independent of $M$ and $\ell_0$.
With $f=\prod f_p$, the preceding calculations give
\[
 f(u^{-1}xu)=\chi_{h_*}(\operatorname{Nrd}u)^{-1}f(x)
       \quad (u\in\widehat R^\times),
 \qquad f(vx)=\alpha(v)f(x)\quad(v\in\widehat{\mathbb Z}^{\times}).
\]
The ramified local quadratic character on units at $p\mid m$ is
$(\cdot/p)$; all other order-unit character factors are trivial.
The scaling character has
\[
 \alpha(-1)=
 \Bigl(\prod_{p^e\Vert N}w_{p^e}\chi_{h_*}(p)^e\Bigr)
             \Bigl(\prod_{p\mid m}(-1/p)\Bigr)(-1)=1.
\]
For the last equality, $m\equiv3\pmod4$ and the first product is
$w(E^{h_*})=+1$.  More explicitly,
$w(E)=-\prod w_{p^e}$ and, since $(h_*,N)=1$,
$w(E^{h_*})=w(E)\chi_{h_*}(-N)$.
The latter twist formula also follows by multiplying the local
unramified-twist factors $\chi_{h_*}(p)^{v_pN}$ at $p\mid N$
and the factors $(-1/p)$ at the newly ramified good primes.

Let $\eta_B$ denote the classical character associated with $\alpha$
in the inverse-idele convention.  The formula
$\eta_B(d)=\prod_p\alpha_p(d)$ for positive integers $d$ prime to
the support is verified in the transformation calculation below.

